\clearpage
\section{PRÉ-DESENVOLVIMENTO}

\subsection{Contextualização do Problema}

Louças sanitárias percorrem uma cadeia composta por fabricação e distribuição, transporte externo, recebimento, movimentação interna e instalação. Como o momento em que o dano se torna visível pode não coincidir com sua origem, a equipe decompôs o fluxo em etapas com atores, operações e riscos distintos. O levantamento inicial combinou mapeamento do fluxo, observação em obra residencial, escuta de atores envolvidos e busca preliminar de anterioridade tecnológica.

\begin{table}[H]
\caption{Etapas da cadeia consideradas na definição do problema}
\label{tab:cadeia}
\centering
\footnotesize
\begin{tabularx}{\textwidth}{L{2.7cm}YYL{2.6cm}}
\toprule
\textbf{Etapa} & \textbf{Operação principal} & \textbf{Sinal observado} & \textbf{Questão de projeto} \\
\midrule
Fabricação e distribuição & Preparar, acondicionar, armazenar e expedir & O setor possui escala industrial e rede distribuída & Como embalagem e unitização afetam o dano latente? \\
Transporte externo & Levar o produto do estoque à obra & O distribuidor associa a esse trecho a maior parte das avarias relatadas & Quais impactos, vibrações, tombamentos e manuseios ocorrem? \\
Recebimento e abertura & Descarregar, conferir e abrir a embalagem & Há relatos de peças que chegam trincadas & Quando o dano fica visível e como conferir com segurança? \\
Movimentação interna & Levar a louça até o ponto de instalação & Há transporte manual quando o carrinho não acessa o local & Quais restrições de acesso precisam ser consideradas? \\
Instalação & Posicionar, fixar e testar & Avarias anteriores podem aparecer nessa etapa & Como evitar retrabalho e exposição a fragmentos? \\
\bottomrule
\end{tabularx}
\fonte{Elaboração da equipe a partir do levantamento da Semana 03 \citep{equipe2026}.}
\end{table}

A visita à obra ofereceu o contexto de recebimento, estoque e movimentação das louças. O distribuidor visitado associou as principais avarias ao transporte entre seu estoque e a obra. Um trabalhador relatou peças que chegam trincadas, quedas durante a descarga e transporte manual quando o acesso impede o uso de carrinho. Esses achados sustentam o recorte inicial, mas ainda não medem frequência, severidade econômica ou contribuição relativa de cada evento logístico \citep{equipe2026}.

A apresentação do projeto registra 22 milhões de peças produzidas por ano no Brasil, 26 unidades fabris de médio e grande porte, presença produtiva em oito estados e aproximadamente sete mil empregos diretos. Os números indicam uma cadeia industrial extensa, mas não fornecem uma taxa de quebra e não permitem estimar perdas sem dados adicionais \citep{equipe2026,anfacer2026}.

\subsection{Contextualização das Dores}

A quebra pode permanecer oculta até o recebimento ou a instalação. Quando isso ocorre, o distribuidor pode absorver reposição e nova logística; a obra enfrenta atraso e retrabalho; e os trabalhadores podem ficar expostos a quedas, esforço adicional e fragmentos cortantes. A formulação do problema separa três questões que deverão ser verificadas: onde o dano surge, quando ele é detectado e quem absorve suas consequências.

\begin{quote}
\itshape
A quebra no transporte entre distribuidor e obra pode permanecer oculta até o recebimento ou a instalação. O evento gera reposição, nova entrega, atraso, desperdício e exposição do trabalhador a peças avariadas ou fragmentos cortantes.
\end{quote}

\subsection{Público Alvo}

O público diretamente afetado inclui distribuidores de louças sanitárias, motoristas e ajudantes responsáveis pelo transporte e pela descarga, responsáveis pelo recebimento, trabalhadores da obra, instaladores, compradores e clientes finais. A definição de quem exerce o papel de comprador, usuário operacional, cliente intermediário e cliente externo ainda será validada no projeto informacional.

\subsection{Legislação e Normas}

\etapafutura{O levantamento das normas aplicáveis ao transporte, à embalagem, à movimentação, à segurança do trabalho e ao descarte será incorporado após a definição da fronteira técnica da solução.}

\pendente{Identificar somente normas e requisitos legais diretamente aplicáveis ao recorte confirmado, registrando versão, item e implicação para o projeto.}

\subsection{Soluções existentes no mercado}

A busca inicial reuniu quatro documentos de patente voltados à proteção de louças ou itens frágeis em transporte. As propostas recorrem a encaixes de papelão, amortecedores moldados, estruturas dobráveis e estruturas corrugadas. O padrão recorrente combina amortecimento, imobilização, resistência e otimização de volume \citep{us2015,jp2007,cn2022,us2019}. O levantamento deverá ser complementado por soluções comercializadas e pela análise de reivindicações, famílias e situação legal das patentes antes de orientar uma concepção.

\subsection{Proposta de Valor}

Como oportunidade de projeto, a equipe formulou a seguinte pergunta:

\begin{quote}
\bfseries
Como reduzir as avarias de louças sanitárias no transporte entre o distribuidor e a obra, diminuindo custos de reposição, atrasos e riscos aos trabalhadores, sem inviabilizar comercialmente o produto ou a operação?
\end{quote}

O enunciado não prescreve forma, material ou mecanismo. Ele admite intervenções no produto, na embalagem, na unitização, no equipamento de movimentação, no processo de conferência ou em uma combinação desses elementos.

\subsection{Custo-Meta}

\etapafutura{O custo-meta será definido depois da identificação do comprador, da estimativa das perdas atuais e da validação da disposição a pagar por uma solução.}

\subsection{Encerramento da fase de Pré-Desenvolvimento}

O pré-desenvolvimento delimitou um problema de cadeia, com fronteira inicial no transporte entre o distribuidor e a obra. A convergência de relatos orienta a investigação, mas não estabelece prevalência ou causalidade. Para avançar, a equipe deverá validar os atores, detalhar o ciclo de vida, registrar as necessidades e obter dados de avaria segmentados por tipo de louça, rota, transportador e etapa de manuseio.
